\section{Associative Memory}

A store you have to search is a poor memory. A good one gives the whole back from a piece, the smell that returns the afternoon, the first three notes that carry the song. This is content-addressable, or associative, memory, recall by what a thing is rather than where it sits, and the mesh is one natively, the computing layer of the model rather than a machine laid over it \cite{potter1992}.

Picture the difference. An ordinary memory is a wall of numbered boxes, and to find a thing you must already know its number. An associative memory keeps no numbers. You hold up a fragment of what you want, and every dock of the store compares that fragment against what it holds, all at once, and the ones that match answer. There is no address to know beforehand, because the content is the address. A partial cue, a few slots of a pattern with the rest left blank, is enough, and the matching pattern returns whole.

The geometry is what makes this both fast and vast. A query spreads from its starting dock as a wave, one ring of neighbors per beat, and the beat a stored match is first touched is the time to recall it. Because the bulk branches exponentially, that wave reaches the entire stored region in a number of beats that grows like the logarithm of its size, where a flat memory of the same width would need a time growing as a root of it, the gap between a glance and a hunt once the store is large. The room grows the same way. The number of distinct patterns that fit within a fixed recall time is the number of docks within that many rings, which on the curved bulk is exponential in the radius and on a flat sheet only a power of it, so against a flat memory of equal reach the curved store holds more from the start and pulls away without bound. There is a second way to pack associations into the same room, binding a key to a value by a reversible combination and bundling many such pairs into a single pattern, each one recovered later by undoing the binding, so one region can carry a whole set of paired associations at once with the capacity again growing with its reach.

\begin{result}[The mesh is an associative memory]
Content-addressable recall is native to the bulk. A partial cue returns the whole stored pattern by a parallel match, the patterns that fit within a fixed recall time grow exponentially with the radius where a flat memory grows polynomially, and a query wave covers the whole region in $O(\log n)$ beats where flat space needs a root of $n$. \cite{pollard2026vibetest}
\end{result}

One thing the memory cannot do on its own, and the limit is as telling as the powers. Hand the bare reversible law a cue corrupted in a fifth of its slots and ask it to clean the noise, and it cannot, the recovered pattern stays near chance. The reason is exact. A reversible law conserves everything, the error included, and cleaning a corrupted memory means throwing the error away, which is forgetting, the one move reversibility forbids. The parallel match recalls perfectly from a true fragment, yet it cannot repair a false one. Repair is a different kind of process, a settling toward the nearest stored pattern that lets the noise drain off, an attractor relaxation of the Hopfield kind, and it needs a small dissipative layer above the reversible base. Add that layer and the same fifth-corrupted cue cleans to better than four parts in five. So the division is clean and worth stating plainly. A memory that never forgets is the reversible base, and a memory that heals itself, that pulls a worn cue back to the thing it almost is, is the work of a dissipative layer, the same lossy drain the arrow and the bath supply everywhere else.
